%%%%%%%%%%%%%%%%%%%%%%%%%%%%%%%%%
%
%       EXERCÍCIO
%
%%%%%%%%%%%%%%%%%%%%%%%%%%%%%%%%%

\ifdefstring{\atividade}{prova}{%
    \renewcommand{\valorquestao}{\ValorQ\ pontos}
}%

\begin{exercicioBanco}[\valorquestao]
Identifique o tipo de amostragem utilizado em cada situação:

\begin{enumerate}[(a)]
    \item Um pesquisador sorteia 30 pacientes de uma lista completa de todos os pacientes atendidos na clínica.
    \item Em um hospital, seleciona-se um paciente a cada 10 registros da agenda.
    \item Os pacientes são separados por faixa etária e, em seguida, são sorteados pacientes dentro de cada faixa.
    \item São sorteadas algumas clínicas de fisioterapia e todos os pacientes dessas clínicas são avaliados.
\end{enumerate}
\end{exercicioBanco}
